\section{Структура и характер цикла в современной экономике}

\subsection{Индикаторы цикла} %\textcolor{red}{Черновик:}
В экономической теории отсутствует единый подход к определению тех статистических показателей, на основании которых можно сделать вывод о наличии циклической динамики, поскольку различные авторы и школы по-разному понимают саму природу цикла, его исходную сферу, механизм развертывания и форму проявления в хозяйственной системе, вследствие чего в одних концепциях решающее значение придается агрегированным данным о совокупной экономической активности, в других – динамике отдельных отраслей, цен, кредита, занятости, инвестиций или структурных сдвигов, а потому вопрос об индикаторах цикла непосредственно связан с более общими теоретическими представлениями о сущности экономической цикличности.

Существующие подходы к выделению индикаторов цикла в зависимости от характера статистических показателей можно разделить на следующие категории:
\begin{itemize}
    \item показатели реального выпуска и производства – реальный ВВП, индекс промышленного производства, производство товаров длительного пользования, объем строительных работ, железнодорожные перевозки, выпуск чугуна и стали; данная группа отражает непосредственную динамику материального производства и совокупного выпуска, поэтому особенно часто используется при датировке фаз цикла;
    \item показатели рынка труда – уровень безработицы, занятость вне сельского хозяйства, совокупные отработанные часы, реальная заработная плата, производительность труда, а также индексы потребительского и предпринимательского доверия; эти показатели позволяют фиксировать изменение спроса на рабочую силу и состояние воспроизводства рабочей силы в различных фазах цикла;
    \item показатели уровня цен и инфляции – общий уровень товарных цен, индекс потребительских цен, дефлятор ВВП, цены производителей, цены на сырье и сельскохозяйственные товары, относительные цены капитальных благ; данная группа особенно значима для теорий, связывающих цикл с колебаниями цен, инфляцией, дефляцией и межсекторными ценовыми диспропорциями;
    \item денежно-кредитные и финансовые показатели – ставка банковского дисконта, ключевая процентная ставка, объем банковского кредита, денежные агрегаты, банковские резервы, клиринговые операции, уровень задолженности, фондовые индексы, спред доходности; эти индикаторы используются в теориях, акцентирующих роль кредита, процента, финансовой хрупкости и долговой нагрузки;
    \item показатели инвестиций и накопления капитала – валовое накопление основного капитала, отношение инвестиций к выпуску, запасы товарно-материальных ценностей, жилищное строительство, ввод новых производственных мощностей, расходы на НИОКР; данная группа позволяет проследить инвестиционную активность и ритм расширенного воспроизводства капитала;
    \item инновационные, технологические и структурные показатели – число патентов, совокупная факторная производительность, временные лаги строительства и ввода мощностей, волны диффузии технологий, смена технологических укладов, норма прибыли и иные структурные характеристики; эти показатели используются прежде всего в теориях инновационных волн, технологических циклов и долгосрочных структурных сдвигов.
\end{itemize}

Наиболее обоснованным является положение о том, что именно показатели динамики промышленного производства выступают наиболее оптимальными индикаторами наличия циклической динамики в экономике, поскольку они, в отличие от многих других статистических рядов, непосредственно отражают движение материального воспроизводства, обладают высокой чувствительностью к изменению хозяйственной конъюнктуры, публикуются с более высокой частотой, чем агрегированные показатели вроде ВВП, и вместе с тем менее опосредованы вторичными перераспределительными, финансовыми или институциональными влияниями, чем показатели цен, кредита или фондового рынка. Именно промышленное производство позволяет сравнительно рано зафиксировать поворотные точки цикла, поскольку в его динамике наиболее отчетливо проявляются колебания выпуска инвестиционных и промежуточных благ, изменение загрузки мощностей, расширение и сокращение спроса на средства производства, то есть те процессы, которые составляют материальную основу промышленного цикла; вследствие этого показатели промышленного производства соединяют в себе достаточную оперативность, высокую чувствительность и прямую связь с воспроизводственными механизмами экономики, что делает их наиболее предпочтительными для выявления циклической динамики.

\subsection{Экономика цикла} %\textcolor{red}{Черновик:}
Различные авторы и теоретические традиции по-разному отвечают на вопрос о том, с какого исторического момента в экономике можно говорить о наличии циклов, в каком именно типе хозяйственной системы они возникают и на каком уровне экономической организации они проявляются наиболее отчетливо, вследствие чего данный вопрос не может быть решен вне общего теоретического подхода к природе самой экономической цикличности.

Теоретические подходы к моменту возникновения цикла в экономике целесообразно классифицировать по следующим историческим периодам:
\begin{itemize}
    \item Древний мир (2000–500 до н. э.) – в рамках данного периода фиксируются лишь наиболее ранние свидетельства повторяющихся хозяйственных колебаний, прежде всего в динамике цен, урожаев и благосостояния, однако речь идет не о теории цикла в строгом смысле, а о практическом наблюдении за повторяемостью отдельных экономических явлений;
    \item Средневековье и раннее Новое время (1260–1700) – в этот период обнаруживаются более устойчивые колебания цен и рыночной конъюнктуры, особенно в сельскохозяйственной сфере, что позволяет ретроспективно говорить о наличии длительных ценовых волн, хотя они еще не обладают характеристиками промышленного делового цикла;
    \item Новое время (1700–1862) – с развитием торговли, кредита и раннего промышленного капитализма циклические колебания приобретают более выраженный и повторяющийся характер; именно к этому периоду многие авторы относят формирование предпосылок современного цикла, прежде всего в Великобритании, Франции и США;
    \item Позднее Новое и Новейшее время (1862 – н. в.) – начиная с работ Жюгляра и последующей теоретической разработки Маркса, Кондратьева, Шумпетера, Митчелла и других исследователей, цикл становится предметом систематического научного анализа, получает четкую фазовую структуру и эмпирическую датировку.
\end{itemize}

При этом, несмотря на наличие более ранних форм хозяйственных колебаний, экономические циклы в собственном, современном смысле слова отчетливо наблюдаются лишь после 1825 года, то есть с момента первого общего кризиса перепроизводства в Великобритании, который обычно рассматривается как первый ``современный'' кризис промышленно-капиталистического типа, соединяющий в себе промышленный, торговый и финансовый компоненты и открывающий эпоху периодически повторяющихся кризисов общего характера.

Точки зрения различных авторов можно классифицировать по следующим типам экономических систем в зависимости от соответствующих им общественно-экономических формаций:
\begin{itemize}
    \item древняя докапиталистическая экономика – здесь могут наблюдаться повторяющиеся колебания цен, урожаев, налоговых поступлений и благосостояния, однако они, как правило, связаны с природно-климатическими, военно-политическими и демографическими факторами и не образуют регулярного цикла перепроизводства;
    \item феодальная экономика – в ее рамках возможны аграрные, ценовые и торговые колебания, однако они носят преимущественно локальный, нерегулярный и внешне обусловленный характер, поскольку сама хозяйственная система еще не основывается на развитом товарном производстве, всеобщей конкуренции и внутренней логике капиталистического накопления;
    \item капиталистическая экономика – именно в рамках капитализма цикл приобретает собственную, внутренне необходимую форму, поскольку возникает из противоречий расширенного воспроизводства капитала, безграничного стремления к накоплению, диспропорциональности развития, кредитной экспансии, изменения нормы прибыли и периодического разрешения накопленных противоречий в форме кризиса. В этом смысле большинство классических и современных теорий – от марксистской и шумпетерианской до кейнсианской, монетарной и австрийской – сходятся в признании того, что именно капиталистическая экономика является основной средой возникновения собственно экономического цикла;
    \item социалистическая экономика – в строгом смысле цикл как закономерное чередование фаз подъема, перепроизводства, кризиса и депрессии не возникает, поскольку при общественной собственности на основные средства производства и планомерной организации воспроизводства отсутствуют анархия производства, конкуренция частных капиталов и всеобщий механизм кризиса реализации, хотя это не исключает возможности диспропорций, ошибок планирования или неравномерности развития.
\end{itemize}

Экономический цикл в собственном смысле слова возникает именно в условиях развитой капиталистической экономики, где производство носит общественный характер, а присвоение – частнокапиталистический, где накопление капитала приобретает всеобщий характер, а рынок становится всеобщей формой связи между производителями. В более ранних формациях отсутствуют необходимые предпосылки промышленного цикла: развернутая система машинного производства, господство наемного труда, кредитная система, развитый внутренний рынок и регулярное перепроизводство капитала. В социалистической же экономике, напротив, устраняются сами базовые причины капиталистического цикла, вследствие чего могут существовать лишь отдельные диспропорции и колебания хозяйственного развития, но не цикл как специфическая форма движения капитала.

В зависимости от того, на каком уровне экономики возникают циклы, различные теоретические подходы можно классифицировать по следующим категориям:
\begin{itemize}
    \item предприятие – на уровне отдельной фирмы или продукта фиксируются колебания продаж, прибыли, инвестиций и жизненного цикла продукции, однако эти изменения слишком зависят от частных обстоятельств, стратегии управления, положения конкретной фирмы на рынке и потому еще не выражают собственно общеэкономический цикл;
    \item отрасль – на отраслевом уровне циклическая динамика проявляется уже гораздо отчетливее, поскольку именно здесь становятся видны волны обновления основного капитала, изменения спроса, отраслевые диспропорции и различия в чувствительности отдельных сфер производства к общему движению экономики;
    \item национальная экономика – это классический уровень наблюдения экономического цикла, поскольку именно здесь агрегируются основные макроэкономические показатели, фиксируются фазы подъема, пика, спада и оживления, а также становится видимым синхронное движение большинства отраслей, занятости, инвестиций и кредита;
    \item региональная экономика – на региональном уровне циклические колебания могут наблюдаться, однако они, как правило, выступают как производные от национального цикла и в значительной мере модифицируются отраслевой структурой региона, степенью его интеграции, специализацией и внешними воздействиями;
    \item мировая экономика – на глобальном уровне можно выявить синхронизацию национальных циклов и длинные волны мирового хозяйственного развития, однако непосредственная циклическая картина здесь усложняется неравномерностью развития стран, различием фаз национальных циклов, действием валютных, политических и геоэкономических факторов.
\end{itemize}

Отчетливо экономические циклы в своем подлинном смысле наблюдаются именно на отраслевом и особенно на национальном уровне: на отраслевом уровне хорошо видны материально-производственные и инвестиционные механизмы цикличности, а на национальном – их агрегированное выражение в движении совокупного выпуска, занятости, инвестиций и кредита. Уровень предприятия является слишком частным, поскольку отдельная фирма может переживать подъем или спад под воздействием индивидуальных факторов, не совпадающих с общеэкономической динамикой. Региональный уровень, в свою очередь, часто отражает лишь специфическую комбинацию отраслевой структуры, территориальной специализации и внешних воздействий, тогда как мировой уровень затрудняет прямое наблюдение цикла вследствие асинхронности национальных хозяйств, неравномерности международного развития и взаимозависимости стран. Поэтому именно отраслевой и национальный уровни выступают основными уровнями, на которых цикл обнаруживается в наиболее явной и научно фиксируемой форме.

\subsection{Исходная сфера цикла} %\textcolor{red}{Черновик:}
Выявление исходной сферы цикла занимает важное место в истории экономической мысли, поскольку различные авторы по-разному определяют тот участок хозяйственной системы, в котором первоначально возникает циклическое движение и откуда оно затем распространяется на остальные сферы воспроизводства, вследствие чего различия между соответствующими теориями касаются не только конкретных механизмов кризиса, но и более общего вопроса о том, где именно следует искать первичный очаг циклической динамики.

Все существующие подходы к исследованию цикла в зависимости от определения его исходной сферы можно свести к двум основным категориям:
\begin{itemize}
    \item сфера производства – к данной группе относятся теории, рассматривающие цикл как результат процессов, возникающих внутри материального производства. Сюда относятся подходы, связывающие начало цикла с диспропорциями между подразделениями общественного воспроизводства, с расширением производства средств производства, с колебаниями инвестиций в основной капитал, с технологическими нововведениями, с волнами предпринимательской активности, а также со строительными и инфраструктурными колебаниями. В этих теориях сфера обмена рассматривается преимущественно как производная или передаточная среда, через которую уже возникшие в производстве противоречия распространяются дальше;
    \item сфера обмена – к данной группе относятся теории, усматривающие исходный импульс цикла в торговле, денежно-кредитной системе или финансовой сфере. Здесь цикл связывается с расширением коммерческого кредита, спекулятивным ростом торговли, накоплением товарных запасов, банковской эмиссией, расхождением между естественной и денежной нормой процента, а также с нарастанием финансовой хрупкости. В рамках данных подходов производство обычно выступает как сфера, на которую уже вторично переносится нарушение, первоначально возникшее в обращении.
\end{itemize}

Однако несмотря на наличие влиятельных концепций, выводящих исходную сферу цикла из обмена, именно промышленное производство выступает исходной сферой цикла, поскольку цикл в собственном смысле слова представляет собой форму движения капиталистического воспроизводства и возникает на почве периодического расширения, перенакопления и кризисного нарушения процесса производства капитала. Именно в промышленности формируются основные материальные предпосылки цикла: диспропорции между подразделениями общественного производства, волны обновления основного капитала, колебания инвестиций и перепроизводство товаров и производственных мощностей, в то время как сфера обмена не является первичным источником цикла, а выполняет производную роль, усиливая, опосредуя и распространяя противоречия, первоначально возникшие в производстве, на торговлю, кредит, денежное обращение и финансовую систему.

\subsection{Фундаментальные причины цикла} %\textcolor{red}{Черновик:}
В экономической теории отсутствует единая общепризнанная трактовка фундаментальных причин экономического цикла, поскольку различные школы и направления по-разному определяют глубинное основание, лежащее в основе циклической динамики.

Все подходы к выделению фундаментальных причин цикла можно разделить по следующим категориям:
\begin{itemize}
    \item противоречия и диспропорции капиталистической экономики – к данной группе относятся теории, объясняющие цикл имманентными свойствами капиталистического способа производства, включая антагонизм между трудом и капиталом, тенденцию нормы прибыли к понижению, анархию рынка, диспропорциональность межотраслевого развития, структурное недопотребление и эндогенную финансовую нестабильность;
    \item нестабильность ожиданий и инвестиций – в рамках этой группы фундаментальная причина цикла связывается с принципиальной неопределенностью инвестиционных решений, изменчивостью предпринимательских ожиданий, психологическими волнами оптимизма и пессимизма, а также со спекулятивными пузырями, нарушающими устойчивость совокупного спроса;
    \item политика и государственное регулирование – данные подходы выводят цикличность из внешнего вмешательства в экономику, прежде всего из ошибок денежно-кредитной политики, искусственного изменения процентной ставки, предвыборного стимулирования экономики, партийной смены макроэкономических приоритетов и, в отдельных концепциях, из целенаправленной геополитической дестабилизации;
    \item закономерности технологического развития – в данной группе цикл объясняется неравномерностью технического прогресса, кластеризацией инноваций, процессами созидательного разрушения, сменой технологических укладов и длинных волн развития, а в ряде теорий – также воздействием случайных технологических шоков.
\end{itemize}

Наиболее последовательным и теоретически обоснованным является положение о том, что фундаментальная причина экономического цикла заключена именно в имманентных противоречиях капиталистической экономики, поскольку только такой подход позволяет рассматривать цикл не как внешнее нарушение изначально устойчивого равновесия и не как результат субъективных ошибок, политических просчетов или изолированных технологических изменений, а как закономерную форму движения капиталистического воспроизводства, тогда как ожидания, инвестиционные колебания, государственная политика и технологические сдвиги лишь изменяют конкретную форму, темп и остроту проявления этих противоречий, но не образуют их первоосновы.

\subsection{Факторы цикла} %\textcolor{red}{Черновик:}
В экономической теории отсутствует единое понимание того, какие именно обстоятельства следует считать ключевыми факторами экономического цикла, поскольку различные школы и направления по-разному определяют те непосредственные условия, под воздействием которых начинается подъем или спад хозяйственной активности, изменяются темпы накопления, нарушаются пропорции воспроизводства и формируются поворотные точки циклической динамики, вследствие чего вопрос о факторах цикла оказывается тесно связанным с более общими представлениями о природе капиталистической экономики как устойчивой либо внутренне нестабильной системы.

Существующие подходы к исследованию факторов цикла целесообразно разделить на три основные категории:
\begin{itemize}
    \item эндогенные – к данной группе относятся подходы, рассматривающие циклические колебания как результат внутренних процессов самой экономической системы, то есть таких изменений, которые порождаются логикой воспроизводства капитала, динамикой инвестиций, кредита, прибыли, занятости, ожиданий и взаимодействием основных макроэкономических переменных;
    \item экзогенные – сюда относятся концепции, связывающие циклические колебания с внешними по отношению к экономической системе воздействиями, такими как природные явления, демографические сдвиги, войны, политические события, случайные технологические шоки или иные импульсы, приходящие извне и нарушающие нормальный ход хозяйственного развития;
    \item смешанные – в рамках данной группы циклическая динамика объясняется сочетанием внешнего импульса и внутренних механизмов его распространения, когда первоначальное воздействие может носить внешний характер, однако его превращение в развернутый цикл становится возможным лишь благодаря внутренним особенностям экономической системы.
\end{itemize}

Все существующие подходы к выделению ключевых факторов цикла можно разделить по следующим категориям:
\begin{itemize}
    \item природные и климатические факторы – включают колебания урожаев, солнечную активность, климатические изменения, а также демографические волны, которые через спрос, цены и строительство могут воздействовать на хозяйственную конъюнктуру;
    \item технологические факторы – охватывают волны инноваций, смену технологических укладов, технологические революции, случайные изменения производительности, а также периодическое обновление основного капитала;
    \item денежно-кредитные факторы – связаны с кредитной экспансией и сжатием, изменением процентной ставки, расхождением между денежной и естественной нормой процента, колебаниями денежного предложения и нарастанием финансовой хрупкости;
    \item инвестиционно-производственные факторы – включают динамику вложений в основной капитал, действие принципа акселерации, лаги в производстве капитальных благ, перепроизводство и чрезмерное расширение производственных мощностей;
    \item доходно-распределительные факторы – связаны с изменением нормы прибыли, соотношения заработной платы и прибыли, уровнем занятости, недопотреблением и воздействием распределительных сдвигов на совокупный спрос и накопление;
    \item психологические факторы – охватывают волны оптимизма и пессимизма, нестабильность ожиданий, “animal spirits”, ошибки предпринимательской оценки будущего спроса, а также стадное поведение и колебания доверия;
    \item политико-институциональные факторы – включают электоральные циклы, войны, изменение режимов регулирования, государственные решения в сфере фискальной и денежно-кредитной политики, а также институциональные перестройки, изменяющие характер хозяйственной динамики.
\end{itemize}

Ключевым фактором экономического цикла выступает именно массовое периодическое обновление основного капитала, поскольку оно непосредственно связывает внутренние противоречия капиталистического воспроизводства с конкретной ритмикой инвестиционного процесса: в ходе расширения производства капиталисты одновременно увеличивают вложения в машины, оборудование, здания и инфраструктуру, что создает волну спроса на средства производства, ускоряет накопление и расширяет масштабы производства, однако по мере завершения инвестиционного бума и насыщения производственного аппарата возникает перенакопление капитала, снижается прибыльность новых вложений, замедляется обновление основного капитала и начинается спад. Именно поэтому периодичность замены крупных элементов производственного аппарата придает циклу его материальную основу и временную определенность, тогда как прочие факторы – денежно-кредитные, психологические, политические или технологические – либо усиливают, либо ослабляют, либо видоизменяют уже складывающуюся циклическую динамику, но не определяют ее ритм столь непосредственно и системно.

\subsection{Период цикла} %\textcolor{red}{Черновик:}
Продолжительность периода экономического цикла, степень его детерминированности и характер тех обстоятельств, которые определяют временные параметры перехода от одной фазы цикла к другой, в экономической теории трактуются далеко не единообразно, поскольку одни исследователи рассматривают цикл как относительно закономерный процесс, связанный с устойчивыми структурными особенностями воспроизводства, другие усматривают в его временной организации лишь вероятностную повторяемость, а третьи вообще отказываются признавать за циклом какую-либо устойчивую длительность, сводя наблюдаемую периодичность к результату взаимодействия случайных воздействий и внутренних механизмов их распространения.

Существующие подходы в зависимости от степени детерминированности цикла в экономике можно разделить на две основные категории:
\begin{itemize}
    \item стохастический – в рамках данного подхода продолжительность цикла не задается жестко и не выводится из какого-либо одного устойчивого механизма, а рассматривается как результат накопления случайных импульсов, внешних шоков, изменений ожиданий или нелинейного распространения разнородных воздействий через экономическую систему;
    \item детерминированный – данный подход исходит из того, что продолжительность цикла обусловлена относительно устойчивыми внутренними параметрами хозяйственного процесса, такими как срок службы основного капитала, длительность производственного процесса, инфраструктурные лаги или институционально заданный политический ритм, вследствие чего цикл обладает определенной, хотя и не абсолютно жесткой, временной закономерностью.
\end{itemize}

Существующие подходы в зависимости от факторов, определяющих продолжительность периода цикла, можно разделить на следующие категории:
\begin{itemize}
    \item материально-производственные факторы – сюда относятся срок физического износа и замены основного капитала, динамика запасов, длительность производственного процесса и структура воспроизводства, то есть те обстоятельства, которые непосредственно связаны с материальной стороной хозяйственной деятельности;
    \item технологические факторы – охватывают волны базовых инноваций, смену технологических укладов, диффузию новых технологий и неравномерность технического прогресса, способные задавать как среднесрочные, так и долгосрочные ритмы развития;
    \item денежно-кредитные факторы – включают кредитную экспансию и сжатие, взаимодействие денежной и естественной нормы процента, накопление финансовой хрупкости и общую динамику банковского кредита, через которые формируется собственный временной ритм подъема и спада;
    \item инфраструктурно-демографические факторы – связаны с волнами жилищного и инфраструктурного строительства, урбанизацией, демографическими сдвигами, изменением числа домохозяйств и иными долгосрочными процессами, влияющими на инвестиционный спрос;
    \item политико-институциональные факторы – включают электоральные циклы, партийно-политические различия, особенности государственного регулирования и институциональной среды, которые могут придавать колебаниям относительно устойчивую продолжительность;
    \item импульсно-стохастические факторы – охватывают случайные внешние импульсы и механизмы их распространения в экономике, когда периодичность не задается заранее, а возникает как статистический результат взаимодействия случайных воздействий с внутренней динамикой системы;
    \item поведенческо-психологические факторы – связаны с изменением ожиданий, стадным поведением, нарастанием и рассеянием спекулятивной эйфории, колебаниями доверия и иными формами коллективной поведенческой динамики, придающими циклу менее устойчивый и более вариативный временной характер.
\end{itemize}

Существующие подходы в зависимости от продолжительности периода цикла можно разделить на следующие категории:
\begin{itemize}
    \item короткие среднесрочные циклы (до 5 лет) – к ним обычно относят циклы запасов и информационных лагов, продолжительность которых составляет приблизительно 3–5 лет и определяется сравнительно быстрыми изменениями в сфере товарного обращения и краткосрочного приспособления производства;
    \item стандартные среднесрочные циклы (7–12 лет) – это наиболее классическая форма экономического цикла, обычно связываемая с инвестициями в основной капитал и его массовым обновлением, что делает данный тип цикла центральным для теории промышленной цикличности;
    \item длинные среднесрочные циклы (15–30 лет) – сюда относятся колебания, связанные главным образом с инфраструктурными инвестициями, строительством, урбанизацией и демографическими волнами, продолжительность которых превышает обычный промышленный цикл;
    \item длинные волны (40 лет и выше) – данная категория охватывает наиболее протяженные ритмы хозяйственного развития, как правило связываемые со сменой технологических укладов, базисных инноваций и долгосрочных структурных преобразований экономики.
\end{itemize}

Наиболее обоснованным является положение о том, что экономический цикл носит относительно детерминированный характер, поскольку его временная структура не является чисто случайной и не сводится к механическому суммированию внешних импульсов, а определяется воспроизводственными закономерностями капиталистической экономики, среди которых решающее значение имеет срок физического и морального износа основного капитала, так как именно периодическое массовое обновление машин, оборудования и других элементов производственного аппарата создает устойчивый ритм инвестиционного спроса, ускоряет накопление в фазе подъема, а затем, по мере завершения инвестиционной волны и насыщения производственных мощностей, ведет к ослаблению вложений и переходу к спаду, вследствие чего наиболее типичной и теоретически оправданной продолжительностью периода цикла следует считать срок около 10 лет, то есть интервал, совпадающий с тем временным горизонтом, который связан с обновлением основной массы производственного капитала.

\subsection{Мультицикличность} %\textcolor{red}{Черновик:}
В экономической теории отсутствует единая точка зрения относительно того, носит ли циклическая динамика экономики одноуровневый характер или же включает несколько различных форм колебательного движения, различающихся по своей продолжительности, происхождению и месту в общей системе воспроизводства, вследствие чего проблема мультицикличности сводится к вопросу о том, следует ли рассматривать наблюдаемую экономическую динамику как проявление одного основного цикла или как результат наложения нескольких циклических процессов.

Все существующие подходы к проблеме мультицикличности можно разделить по следующим категориям:
\begin{itemize}
    \item существование единственного цикла при отсутствии остальных – в рамках данной позиции любые наблюдаемые колебания сводятся к одному основному механизму, тогда как остальные предполагаемые циклы рассматриваются либо как статистическая иллюзия, либо как производные формы того же самого процесса;
    \item существование доминирующего цикла при зависимом положении остальных – данный подход признает наличие нескольких циклов, однако исходит из того, что один из них задает общий ритм и основную логику экономического движения, в то время как остальные накладываются на него, модифицируют отдельные фазы или зависят от его более широкого контекста;
    \item существование нескольких равнозначных независимых циклов – здесь предполагается, что различные циклы обладают собственными причинами, собственной продолжительностью и собственной сферой действия, сосуществуя параллельно без жесткой иерархии и без обязательной причинной зависимости друг от друга;
    \item существование нескольких равнозначных взаимозависимых циклов – в рамках этой позиции циклы рассматриваются как различные, но взаимно влияющие друг на друга процессы, которые могут усиливаться, ослабляться, синхронизироваться или вступать в резонанс, не образуя при этом строгого отношения подчинения одного другому.
\end{itemize}

Однако наиболее обоснованным с теоретической и практической точек зрения является существование единственного доминирующего промышленного цикла, тогда как остальные циклические колебания занимают подчиненное положение, поскольку именно промышленный цикл наиболее непосредственно выражает логику капиталистического воспроизводства, связанную с накоплением капитала, массовым обновлением основного капитала, расширением и последующим сжатием инвестиций, тогда как краткосрочные колебания запасов, кредитные и финансовые волны, строительные ритмы и прочие колебания либо обслуживают этот основной процесс, либо модифицируют условия его протекания, либо отражают его на иных временных уровнях. Иначе говоря, прочие циклы не обладают тем же уровнем системообразующего значения, что и промышленный цикл, поскольку они либо уже по охвату, либо вторичны по отношению к движению производственного капитала, вследствие чего именно промышленный цикл следует рассматривать как главный, а остальные – как зависимые и менее значимые по своему теоретическому и воспроизводственному содержанию.

\subsection{Фазы цикла} %\textcolor{red}{Черновик:}
В экономической теории отсутствует единый подход к определению фаз экономического цикла, поскольку различные авторы по-разному решают вопрос как о количестве фаз, так и об их содержании, наименовании и последовательности, что обусловлено различием исходных теоретических установок, а именно тем, что именно считается нормальным состоянием экономики, какой момент принимается за начало цикла и какие процессы признаются решающими в переходе от одной стадии циклического движения к другой.

Существующие подходы к выделению фаз цикла можно разделить на следующие категории:
\begin{itemize}
    \item модели с двумя фазами – сводят цикл к наиболее общей бинарной схеме, как правило, в виде подъема и спада, либо бума и депрессии, что обеспечивает простоту анализа, но делает описание внутренней структуры цикла слишком укрупненным;
    \item модели с тремя фазами – выделяют, помимо подъема и спада, особую переломную стадию кризиса или ликвидации, благодаря чему цикл получает более содержательную структуру, однако переходные состояния все еще фиксируются недостаточно полно;
    \item модели с четырьмя фазами – представляют наиболее распространенный и уравновешенный вариант периодизации, поскольку позволяют выделить не только подъем и спад, но и промежуточные состояния депрессии и оживления либо, в иной терминологии, пик и дно, тем самым формируя целостную картину движения цикла;
    \item модели с пятью фазами – обычно связаны с более детализированным анализом финансовой и психологической динамики, когда внутри общего цикла дополнительно выделяются стадии эйфории, фиксации прибыли, паники или поворотной точки;
    \item модели с восемью и более фазами – дают максимально подробное описание циклического процесса, расчленяя его на множество последовательных состояний, что повышает аналитическую детализацию, но одновременно затрудняет обобщение и практическое использование такой схемы.
\end{itemize}

Наиболее обоснованной и методологически оптимальной является четырехфазовая модель цикла с последовательностью кризис → депрессия → оживление → подъем, поскольку она, с одной стороны, избегает чрезмерного упрощения, характерного для двух- и трехфазных схем, а с другой – не дробит циклический процесс на избыточное количество промежуточных состояний, как это происходит в пятифазных и многофазных моделях, вследствие чего именно она позволяет наиболее последовательно отразить внутреннюю логику промышленного цикла: кризис выступает как момент разрыва прежнего воспроизводственного равновесия и начало нового цикла, депрессия – как стадия застоя, обесценения капитала и подготовки условий для нового расширения, оживление – как переход к восстановлению производства и инвестиций, а подъем – как фаза ускоренного накопления, в рамках которой постепенно созревают предпосылки нового кризиса; именно поэтому четырехфазная модель не только сохраняет необходимую аналитическую ясность, но и наиболее полно выражает логику капиталистического воспроизводства.